% TEMPLATE-MILC-v3.tex - plantilla maestra UNICA de las guias MILC v3.
%
% La usan las tres familias de guias, cada una con su driver:
%   · Grados 6-11  -> scripts/build-guias-g11.py
%   · Bebras       -> scripts/build-guias-bebras.py
%   · Semillero    -> scripts/build-guias-semillero.py
%
% Antes habia tres copias casi identicas (~400 lineas iguales, 6 distintas):
% cada mejora habia que aplicarla tres veces, y en la practica no se hacia
% (el motor de recursos vivia solo en dos de ellas). Lo que varia entre
% familias esta parametrizado en 6 placeholders que cada driver llena
% (se nombran aqui SIN su sintaxis real: el builder reemplaza por texto plano,
% tambien dentro de los comentarios, y un valor multilinea rompe el preambulo):
%
%   HEADER_IZQ        encabezado izquierdo de pagina
%   HEADER_DER        encabezado derecho de pagina
%   PORTADA_ETIQUETA  rotulo sobre el numero grande (GRADO/MOMENTO/MODULO)
%   PORTADA_SERIE     linea de serie de la portada
%   PORTADA_ID        identificador de la guia en la portada
%   PORTADA_CIRCULO   el circulito de grado (°), o vacio si no aplica
%
% El resto de claves las documenta content/guias/_SCHEMA.md. El script valida
% que no queden placeholders sin reemplazar antes de escribir el archivo final.

\documentclass[11pt,letterpaper]{article}
\usepackage[margin=1.75cm]{geometry}
\usepackage[table]{xcolor}
\usepackage{fontspec}
\usepackage[spanish]{babel}
\usepackage{tikz}
% Librerías TikZ para los diagramas de las guías (bloque `recursos:`):
%  · babel  → babel en español vuelve activos caracteres como «>», y TikZ los usa
%             en sus opciones (>=stealth, ->). Sin esta librería, cualquier
%             diagrama con flechas revienta con «\language@active@arg> has an
%             extra }». Debe ir ANTES de que se dibuje nada.
%  · positioning        → sintaxis «below left=2pt and 3pt».
%  · decorations.pathreplacing → llaves ({brace}) para señalar rangos.
\usetikzlibrary{babel,positioning,decorations.pathreplacing}
\usepackage{graphicx}
% Circuitos: el trabajo de electrónica se hace hoy con micro:bit y sus sensores
% integrados (MakeCode, sin cableado), así que no se necesitan esquemáticos.
% Si algún día entran montajes con protoboard, basta descomentar esta línea:
% los diagramas .tex/.tikz declarados en `recursos:` ya se incrustan solos.
% \usepackage{circuitikz}
\usepackage[most]{tcolorbox}
\usepackage{tabularx}
\usepackage{array}
\usepackage{fancyhdr}
\usepackage{titlesec}
\usepackage{ragged2e}
\usepackage{enumitem}
\usepackage{microtype}

% Helvetica Neue no trae la flecha U+2192, asi que cada «→» escrito literalmente
% en el YAML se compilaba a NADA, en silencio (462 flechas en 61 guias: rutas del
% tipo «paleta → tipografia → lockup» salian rotas). Mapeamos el caracter a su
% version matematica, que si tiene glifo. Asi el autor puede seguir escribiendo
% «→» comodamente en el YAML.
\usepackage{newunicodechar}
\newunicodechar{→}{\ensuremath{\rightarrow}}

\setmainfont{Helvetica Neue}
\setsansfont{Helvetica Neue}
\setlength{\parindent}{0pt}
\setlength{\parskip}{6pt}
\hyphenpenalty=200
\exhyphenpenalty=200
\pretolerance=400
\tolerance=2000
\setlength{\emergencystretch}{1em}
\renewcommand{\arraystretch}{1.25}
\setlist{leftmargin=5mm,itemsep=1mm,topsep=1mm}
\newcolumntype{Y}{>{\RaggedRight\arraybackslash}X}

\definecolor{milcMagenta}{HTML}{E6007E}
\definecolor{milcVino}{HTML}{5A0038}
\definecolor{milcTurquesa}{HTML}{0093A5}
\definecolor{milcVerde}{HTML}{58A923}
\definecolor{milcMostaza}{HTML}{E5B400}
\definecolor{milcOcre}{HTML}{B86600}
\definecolor{milcNegro}{HTML}{241816}
\definecolor{milcGris}{HTML}{F7F7F4}
\definecolor{milcLinea}{HTML}{E8E2DD}
\definecolor{milcGradePrimary}{HTML}{7C3AED}
\definecolor{milcGradeDark}{HTML}{4C1D95}
\definecolor{milcGradeSoft}{HTML}{EDE9FE}
\definecolor{milcGradeAccent}{HTML}{CFE7FF}
\definecolor{milcGradeText}{HTML}{FFFFFF}
\color{milcNegro}

\pagestyle{fancy}
\fancyhf{}
\lhead{\small Tecnología e Informática · Grado 9}
\rhead{\small Guía 21 · MILC v3}
\cfoot{\small\thepage}
\renewcommand{\headrulewidth}{0pt}

\newtcolorbox{titlebox}[1]{
  enhanced,colback=#1,colframe=#1,arc=7mm,boxrule=0pt,
  left=7mm,right=7mm,top=3mm,bottom=3mm,
  before skip=8pt,after skip=8pt,
  fontupper=\sffamily\bfseries\Large\color{white}
}

\newtcolorbox{softbox}[3]{
  enhanced,colback=#2,colframe=#1,borderline west={4pt}{0pt}{#1},
  arc=2mm,boxrule=.4pt,left=4mm,right=4mm,top=3mm,bottom=3mm,
  before skip=6pt,after skip=8pt,title={#3},coltitle=white,
  colbacktitle=#1,fonttitle=\sffamily\bfseries,fontupper=\RaggedRight,
  attach boxed title to top left={xshift=3mm,yshift=-2mm},
  boxed title style={arc=2mm, boxrule=0pt}
}

\newtcolorbox{drawbox}[2]{
  enhanced,colback=white,colframe=#1,arc=4mm,boxrule=1pt,
  left=5mm,right=5mm,top=4mm,bottom=4mm,before skip=8pt,after skip=8pt,
  title={#2},coltitle=white,colbacktitle=#1,fonttitle=\sffamily\bfseries,
  fontupper=\RaggedRight,
  attach boxed title to top left={xshift=4mm,yshift=-2mm},
  boxed title style={arc=3mm, boxrule=0pt}
}

\newtcolorbox{infoband}[1]{
  enhanced,colback=#1!8,colframe=#1!50,arc=2mm,boxrule=.5pt,
  left=4mm,right=4mm,top=2mm,bottom=2mm,before skip=4pt,after skip=4pt,
  fontupper=\small\RaggedRight
}

% Puente narrativo entre secciones (contrato editorial Conéctate).
% Da continuidad: "Acabas de X. Ahora vas a Y porque Z."
% Visualmente: banda izquierda en color del grado, texto en itálica pequeña.
\newtcolorbox{puentebox}{
  enhanced,colback=milcGradeSoft,colframe=milcGradeDark,
  borderline west={3pt}{0pt}{milcGradeDark},
  arc=0mm,boxrule=0pt,left=5mm,right=4mm,top=2.5mm,bottom=2.5mm,
  before skip=4pt,after skip=8pt,
  fontupper=\itshape\small\color{milcNegro}\RaggedRight
}
% La flecha va en modo matematico a proposito: Helvetica Neue NO tiene el glifo
% U+2192, asi que \textrightarrow se compilaba a NADA («Missing character: There
% is no → in font Helvetica Neue») y el puente salia sin flecha. En modo
% matematico la flecha viene de la fuente matematica y si se dibuja.
\newcommand{\puente}[1]{%
  \begin{puentebox}%
    \textcolor{milcGradeDark}{$\rightarrow$}\hspace{2mm}#1%
  \end{puentebox}%
}

\newcommand{\linead}[1]{\noindent\color{milcLinea}\rule{#1}{0.5pt}\color{milcNegro}}
\newcommand{\respuesta}[1]{\par\vspace{2mm}\foreach \n in {1,...,#1}{\noindent\color{milcLinea}\rule{\linewidth}{0.5pt}\par\vspace{5mm}}\color{milcNegro}}
\newcommand{\checkbox}{\textcolor{milcGradeDark}{\fbox{\phantom{x}}}\hspace{5pt}}

\newcommand{\phasecard}[4]{
\begin{tcolorbox}[enhanced,colback=#3,colframe=#2,arc=3mm,boxrule=.5pt,height=3.05cm,valign=center,halign=center,left=1.5mm,right=1.5mm,top=2mm,bottom=2mm]
{\sffamily\bfseries\fontsize{22}{24}\selectfont\textcolor{#2}{#1}}\par
{\sffamily\bfseries\small #4}
\end{tcolorbox}
}

% ── Recursos visuales de la guía ──────────────────────────────────────
% Declarados en el bloque `recursos:` del YAML y emitidos por el builder.
% \guiaFigura   → imagen rasterizada o PDF (foto, carta celeste, montaje).
% \guiaDiagrama → fuente TikZ/circuitikz (.tex) incrustada como vector:
%                 es la vía para circuitos y esquemáticos editables.
% #1 = ruta relativa al .tex · #2 = pie (puede ir vacío)
\newcommand{\guiaFigura}[2]{%
\begin{center}
\includegraphics[width=0.88\linewidth,height=0.38\textheight,keepaspectratio]{#1}\\[1.5mm]
{\footnotesize\itshape\textcolor{milcNegro!75}{#2}}
\end{center}
}

\newcommand{\guiaDiagrama}[2]{%
\begin{center}
\input{#1}\\[1.5mm]
{\footnotesize\itshape\textcolor{milcNegro!75}{#2}}
\end{center}
}

\begin{document}
\thispagestyle{empty}
\begin{tikzpicture}[remember picture,overlay]
  \fill[white] (current page.south west) rectangle (current page.north east);
  \path[fill=milcGradePrimary,rounded corners=16mm]
    ([xshift=.55cm,yshift=-.55cm]current page.north west) rectangle
    ([xshift=-.55cm,yshift=3.65cm]current page.south east);

  \node[anchor=north west,text=milcGradeText] at ([xshift=2.0cm,yshift=-1.85cm]current page.north west)
    {\sffamily\bfseries\fontsize{14}{16}\selectfont GRADO};
  \node[anchor=north west,text=milcGradeText,opacity=.82] at ([xshift=2.0cm,yshift=-2.75cm]current page.north west)
    {\sffamily\bfseries\fontsize{9}{11}\selectfont SERIE GUÍAS · TECNOLOGÍA E INFORMÁTICA};

  \begin{scope}[shift={([xshift=-3.05cm,yshift=-2.55cm]current page.north east)}]
    \fill[white,opacity=.80] (-.62,-.52) rectangle (.62,.52);
    \draw[milcGradeText,line width=1pt] (-.62,-.52) rectangle (.62,.52);
    \fill[milcGradeDark] (-.42,-.38) rectangle (-.22,.06);
    \fill[milcGradeAccent] (-.10,-.38) rectangle (.10,.24);
    \fill[milcGradePrimary!60!black] (.22,-.38) rectangle (.42,.38);
  \end{scope}

  \node[anchor=west,text=milcGradeText] at ([xshift=1.9cm,yshift=-12.85cm]current page.north west)
    {\sffamily\bfseries\fontsize{124}{124}\selectfont 9};
  % El circulito de grado (°) solo aplica a las guias de grado («11°»); en
  % Bebras y Semillero el numero grande es un momento/modulo, no un grado.
  \draw[milcGradeText,line width=1.7pt]
    ([xshift=4.52cm,yshift=-11.68cm]current page.north west) circle (.13cm);

  \node[anchor=west,text=milcGradeText,text width=8.7cm,align=left] at ([xshift=10.35cm,yshift=-11.6cm]current page.north west)
    {
     {\sffamily\bfseries\fontsize{19}{22}\selectfont ¿Qué es un dato? ---\\recolección\\con propósito}\\[4mm]
     {\sffamily\bfseries\fontsize{23}{26}\selectfont Guía 21-9-TIC}\\[2mm]
     {\sffamily\fontsize{13}{16}\selectfont Período 3 · Datos --- del registro al insight}\\[5mm]
     {\sffamily\bfseries\fontsize{11}{13}\selectfont Institución Educativa Sor María Juliana}\\[1mm]
     {\sffamily\fontsize{10}{12}\selectfont Autor: Dr. Álvaro Cárdenas Orozco}
    };

  \path[fill=milcGradeSoft,rounded corners=7mm]
    ([xshift=1.35cm,yshift=.85cm]current page.south west) rectangle
    ([xshift=-1.35cm,yshift=4.25cm]current page.south east);
  \draw[milcGradeDark,line width=.65pt,rounded corners=7mm]
    ([xshift=1.35cm,yshift=.85cm]current page.south west) rectangle
    ([xshift=-1.35cm,yshift=4.25cm]current page.south east);
  \node[anchor=center,text=milcNegro,text width=17.0cm,align=center] at ([yshift=2.55cm]current page.south)
    {\sffamily\fontsize{10}{12}\selectfont Metodología MILC · Apertura ancestral · Pensamiento computacional · Triángulo Dussel-Estoicismo-Floridi\\
    \textcolor{milcGradeDark}{\bfseries Producto:} 1 hoja de cálculo con 20 datos reales recolectados por el estudiante + una pregunta concreta que esos datos van a responder};
\end{tikzpicture}
\mbox{}
\newpage

\section*{Apertura: ¿Qué es un dato? --- recolección con propósito}

\begin{softbox}{milcGradeDark}{milcGradeSoft}{Saber ancestral}
En las \textbf{fincas cafeteras del norte del Valle del Cauca} (Argelia, Ulloa, Cartago, Versalles) y en las \textbf{veredas Páez del Cauca}, el \emph{conteo de cosecha} nunca fue vanidad estadística. La abuela contaba en su libreta los \emph{racimos por mata} de café, las \emph{mazorcas por surco}, las \emph{mantas terminadas por luna}, los \emph{días que llovió}. Los \textbf{taitas Páez (Nasa)} llevaban en piedrecitas y nudos el registro del año productivo de la comunidad. Cada cifra tenía propósito: saber cuánto sembrar el próximo ciclo, a quién pedir ayuda, qué reservar para el invierno, qué compartir con la asamblea. \textbf{Sin propósito, contar es desperdicio.} La estadística académica olvidó esa lección; el campesino del Valle y el comunero indígena del Cauca nunca la olvidaron.
\end{softbox}

\begin{softbox}{milcTurquesa}{milcGris}{Saber contemporáneo}
Hoy un \textbf{dato} es cualquier información cuantificable o categorizable que puedas registrar de forma estructurada: edad, hora, lugar, color, número, sí/no. Pero ``recolectar datos'' sin pregunta clara es lo más común de la era digital ---  apps que registran todo, hojas de cálculo que crecen sin propósito. La phronesis del dato empieza por ahí: \emph{¿qué pregunta voy a responder?}. Sin pregunta, no hay análisis posible; con pregunta, cada dato cuenta.
\end{softbox}

\begin{softbox}{milcMostaza}{milcGris}{Pregunta-puente}
\textit{¿Qué datos recogía tu abuela en su libreta de cosecha y para qué? ¿Y qué datos recoges tú a diario (en redes, en apps, en notas) sin tener pregunta clara que esos datos vayan a responder?}
\end{softbox}

\begin{softbox}{milcVerde}{milcGris}{Saber y saber hacer}
Al terminar podrás: (1) \textbf{identificar} la diferencia entre dato, información y ruido; (2) \textbf{analizar} 3 prácticas de recolección de datos reales (encuesta, observación, registro propio); (3) \textbf{aplicar} la regla ``pregunta antes que dato'' a un caso propio; (4) \textbf{crear} tu primera tabla de 20 datos reales recolectados con pregunta declarada.
\end{softbox}



\small
\begin{tabularx}{\textwidth}{>{\bfseries}p{3.0cm}Y}
\rowcolor{milcGradePrimary!12}Autor & Dr. Álvaro Cárdenas Orozco\\
DBA & Tratamiento y visualización honesta de datos para la toma de decisiones (MEN, T\&I)\\
Referentes & Ética del dato · Visualización honesta · Pensamiento computacional\\
Producto final & 1 hoja de cálculo con 20 datos reales recolectados por el estudiante + una pregunta concreta que esos datos van a responder\\
\end{tabularx}
\normalsize

\puente{Antes de empezar, mira el viaje completo del Periodo 3. En 10 sesiones vas a recorrer el ciclo del dato completo: recolección $\rightarrow$ tablas $\rightarrow$ fórmulas $\rightarrow$ filtros $\rightarrow$ tablas dinámicas $\rightarrow$ gráficos $\rightarrow$ limpieza $\rightarrow$ honestidad visual $\rightarrow$ insights $\rightarrow$ sustentación de un mini-estudio. Hoy empiezas por la base: \emph{¿qué dato vale recolectar y para qué?}.}

\begin{titlebox}{milcGradeDark}
Ruta MILC: así vamos a aprender
\end{titlebox}
\begin{tabularx}{\textwidth}{>{\centering\arraybackslash}X>{\centering\arraybackslash}X>{\centering\arraybackslash}X>{\centering\arraybackslash}X}
\phasecard{1}{milcVerde}{milcGris}{Escucha\\[-1mm]\small Observo, escucho y pregunto.} &
\phasecard{2}{milcTurquesa}{milcGris}{Sistematización\\[-1mm]\small Pensamiento computacional.} &
\phasecard{3}{milcMagenta}{milcGris}{Praxis\\[-1mm]\small Construyo y aplico.} &
\phasecard{4}{milcVino}{milcGris}{Evaluación\\[-1mm]\small Reflexión liberadora.}
\end{tabularx}


\puente{Antes de teorizar, vas a auditar tu propia relación con los datos: cuántos datos tuyos están recogiendo ya las apps, cuántos recoges tú, con qué propósito. Esa auditoría te muestra que vives rodeado de datos sin pregunta.}

\begin{titlebox}{milcVerde}
Fase 1 · Escucha
\end{titlebox}

\begin{softbox}{milcVerde}{milcGris}{Escena para escuchar}
\textbf{Actividad 1 · IDENTIFICA --- Tu propia huella de datos} (15 min · individual). Abre tu celular y haz inventario: (a) ¿qué \emph{apps} recogen datos tuyos (ubicación, tiempo de uso, contactos)?; (b) ¿qué \emph{datos cotidianos} registras tú (notas en cuaderno, lista de gastos, plan semanal)?; (c) ¿qué \emph{datos te falta} registrar de tu propia vida que mejorarían una decisión?. Anota cada lista con honestidad.
\end{softbox}

\checkbox Identifiqué 5+ apps que recogen datos míos sin que yo decida.\\
\checkbox Identifiqué 3 prácticas de registro propio (lista de compras, notas, agenda).\\
\checkbox Reconocí al menos 1 dato que NO registro y me serviría para mejorar una decisión.

\begin{infoband}{milcVerde}


\textbf{Lo que va al cuaderno (Actividad 1 · IDENTIFICA).} Título: «Actividad 1 --- Mi huella de datos». \textbf{Formato:} 3 listas (Apps que me miden $\mid$ Datos que registro yo $\mid$ Datos que me servirían y no registro) con 5+ ítems cada una. \textbf{Sabes que terminaste cuando} reconoces honestamente que vives rodeado de datos pero recoges pocos con propósito propio.
\end{infoband}

\puente{Ya viste tu propia situación. Ahora vamos a entender los \textbf{4 elementos} de toda recolección de datos seria: pregunta acotada, instrumento honesto, muestra adecuada, registro estructurado. Sin alguno, los datos recolectados son ruido.}

\begin{titlebox}{milcTurquesa}
Fase 2 · Sistematización con pensamiento computacional
\end{titlebox}

\begin{softbox}{milcTurquesa}{milcGris}{Cuatro pilares aplicados al tema}
Toda recolección de datos seria cumple 4 condiciones. Vamos a descomponerlas con los pilares del pensamiento computacional para que tu primera tabla no sea improvisada.
\end{softbox}

\small
\begin{tabularx}{\textwidth}{>{\bfseries\RaggedRight\arraybackslash}p{3.6cm}Y}
\rowcolor{milcTurquesa!90}\color{white}Pilar & \color{white}\bfseries Aplicación al tema\\
1. Descomposición & Una recolección honesta se descompone en 4 elementos: (1) \textbf{pregunta acotada} (qué quieres responder con esos datos, en una frase corta), (2) \textbf{instrumento honesto} (encuesta, observación o registro que no induzca respuesta), (3) \textbf{muestra adecuada} (a quién mides y cuántos: ni 3 amigos, ni todo el colegio sin saber por qué), (4) \textbf{registro estructurado} (tabla con columnas tipadas, no notas sueltas).\\
2. Reconocimiento de patrones & Toda recolección sigue el patrón \textbf{``pregunta primero, dato después''}: cuando alguien dice ``vamos a recoger datos del curso'' sin pregunta clara, lo que se recoge es ruido. Con pregunta acotada (``¿cuánto tiempo le dedican a leer fuera del colegio?''), cada dato encuentra su lugar en la tabla y su uso en el análisis.\\
3. Abstracción & Lo esencial cabe en una pregunta: \emph{¿esta recolección la podré usar la próxima semana para decidir algo concreto?}. Si la respuesta es no, todavía no es recolección --- es curiosidad sin método.\\
4. Algoritmo conceptual & Pasos para recolectar bien: (1) escribe la pregunta en 1 frase; (2) decide instrumento (encuesta/observación/registro); (3) elige muestra (a quién y cuántos); (4) diseña tabla con columnas tipadas; (5) recoge sin trampa; (6) revisa que cada fila esté completa antes de cerrar.\\
\end{tabularx}
\normalsize

\begin{drawbox}{milcTurquesa}{Anatomía de una recolección honesta}
\textbf{Pregunta:} 1 frase clara que responde la recolección. \\[1mm] \textbf{Instrumento:} cómo vas a recoger (encuesta, observación, registro propio). \\[1mm] \textbf{Muestra:} a quién mides (cuántos, qué grupo). \\[1mm] \textbf{Tabla:} columnas tipadas (texto, número, fecha, sí/no). \\[1mm] \textbf{Registro:} llenado completo, sin filas vacías. \\[1mm] \textbf{Uso previsto:} qué vas a decidir con esos datos.
\end{drawbox}

\begin{softbox}{milcMostaza}{milcGris}{Errores comunes y su corrección}
(1) Recolectar sin pregunta: el clásico ``vamos a hacer una encuesta'' sin propósito. (2) Pregunta demasiado vaga (``cómo va el curso''): da datos basura. (3) Muestra solo de amigos cercanos: sesgo de origen. (4) Mezclar tipos en una misma columna (números y texto): rompe el análisis posterior. (5) Olvidar el uso previsto: recolectas y nunca decides nada.
\end{softbox}

\begin{infoband}{milcTurquesa}


\textbf{Lo que va al cuaderno (Actividad 2 · ANALIZA).} Título: «Actividad 2 --- 4 elementos de una recolección honesta». \textbf{Formato:} 4 fichas, una por elemento, cada una con: definición en tus palabras + ejemplo bueno + ejemplo malo (uno que conozcas, incluso de tus propias tareas anteriores). \textbf{Sabes que terminaste cuando} reconoces al menos UN error que has cometido tú en recolecciones anteriores.
\end{infoband}

\puente{Tienes el método. Toca producir: una hoja con 20 datos reales recolectados, instrumento documentado, pregunta acotada al inicio. Es la base de todo el periodo.}

\begin{titlebox}{milcMagenta}
Fase 3 · Praxis con pensamiento computacional
\end{titlebox}

\begin{softbox}{milcMagenta}{milcGris}{Construyendo el producto con los cuatro pilares}
\textbf{Actividad 3 · CREA --- Tu primera tabla con 20 datos reales} (30 min · individual o pareja). Hora de recolectar. Elige UNA pregunta acotada y recoge 20 filas de datos reales en una hoja de cálculo (Google Sheets, Excel o cuaderno digital). Cada fila debe responder a la pregunta.
\end{softbox}

\small
\begin{tabularx}{\textwidth}{>{\bfseries\RaggedRight\arraybackslash}p{3.6cm}Y}
\rowcolor{milcMagenta!90}\color{white}Pilar & \color{white}\bfseries Aplicación al producto\\
Descomposición & Tu recolección se descompone en 5 pasos: (1) escribir la pregunta en 1 frase, (2) decidir instrumento (encuesta breve a 20 compañeros, observación de 20 hechos, registro propio de 20 días), (3) crear tabla con 3-5 columnas tipadas, (4) recoger los 20 datos sin saltarse columnas, (5) verificar que cada fila esté completa.\\
Patrones & Sigue el patrón \textbf{``ejemplo concreto vence a categoría abstracta''}: tu pregunta debe ser específica. Mal: ``¿cómo usan el celular?''. Bien: ``¿cuántos minutos diarios pasan mis compañeros en TikTok los lunes?''. La especificidad permite responder.\\
Abstracción & Cada columna expresa un solo tipo de dato. Si una columna mezcla números con texto, divídela en dos columnas distintas. La pureza tipográfica es base del análisis.\\
Algoritmo de armado & Algoritmo: (1) elige tema; (2) escribe pregunta acotada; (3) diseña 3-5 columnas con su tipo; (4) recoge 20 filas reales; (5) revisa que ninguna fila esté incompleta; (6) guarda el archivo con nombre claro.\\
\end{tabularx}
\normalsize

\begin{drawbox}{milcMagenta}{Checklist antes de cerrar la tabla}
\checkbox Pregunta acotada en 1 frase escrita arriba de la tabla. \\[1mm] \checkbox 3-5 columnas con tipos declarados (texto, número, fecha, sí/no). \\[1mm] \checkbox 20 filas reales recolectadas (no inventadas). \\[1mm] \checkbox Cero filas con celdas vacías. \\[1mm] \checkbox Archivo guardado con nombre que incluye tu nombre + pregunta corta. \\[1mm] \checkbox Puedes explicar qué vas a decidir con esos 20 datos.
\end{drawbox}

\begin{softbox}{milcMagenta}{milcGris}{Plantilla de trabajo}
\textbf{Pregunta acotada.} \textit{1 frase corta y específica.} \\[1mm] \textbf{Instrumento.} \textit{Cómo recogí los datos: encuesta, observación o registro propio.} \\[1mm] \textbf{Muestra.} \textit{A quién encuesté o qué observé.} \\[1mm] \textbf{Tabla.} \textit{3-5 columnas con tipo + 20 filas reales.} \\[1mm] \textbf{Uso previsto.} \textit{Qué voy a decidir con estos datos.}
\end{softbox}

\begin{infoband}{milcMagenta}


\textbf{Lo que va al cuaderno (Actividad 3 · CREA).} Título: «Actividad 3 --- Mi primera tabla de 20 datos». \textbf{Formato:} captura impresa o link de la hoja + pregunta + instrumento + uso previsto + reflexión de 3 renglones. \textbf{Extensión:} 1 página de cuaderno + 1 archivo digital. \textbf{Sabes que terminaste cuando} tu hoja tiene 20 filas reales y sirve para responder UNA pregunta concreta.
\end{infoband}

\puente{Lo que sigue resume \emph{qué} entregas hoy. El criterio principal: que cada uno de los 20 datos sirva a la pregunta que te planteaste, y que la pregunta sea concreta (no ``cómo es el curso'').}

\begin{titlebox}{milcMostaza}
Producto de la sesión
\end{titlebox}
\textbf{Tabla de 20 datos reales recolectados con pregunta acotada y uso previsto}.
\begin{softbox}{milcMostaza}{milcGris}{Criterios de calidad}
(1) Pregunta acotada en 1 frase clara. (2) 20 filas de datos reales (no inventadas). (3) 3-5 columnas tipadas correctamente. (4) Cero filas incompletas. (5) Uso previsto declarado. (6) El archivo se puede abrir y leer sin explicación oral.
\end{softbox}

\puente{Antes de cerrar, mira tu recolección desde las cinco dimensiones humanas. Decidir qué datos recoger es decisión política sutil: quién aparece, quién queda fuera, qué cuenta y qué no.}

\begin{titlebox}{milcVino}
Fase 4 · Evaluación liberadora — cinco dimensiones
\end{titlebox}

\begin{softbox}{milcVino}{milcGris}{Sentido de la evaluación}
Evaluar no es solo revisar una nota. Es reconocer qué aprendí, qué cambió en mí, qué emociones manejé, cómo me relaciono con la ciudad y la comunidad, y qué le dejaré a quien viene detrás.
\end{softbox}

\textbf{Mi reflexión liberadora en cinco dimensiones.} Completa cada frase iniciada:

\small
\begin{tabularx}{\textwidth}{>{\bfseries\RaggedRight\arraybackslash}p{3.4cm}Y>{\RaggedRight\arraybackslash}p{4.6cm}}
\rowcolor{milcVino}\color{white}Dimensión & \color{white}\bfseries Mi respuesta (frame starter) & \color{white}\bfseries Algo que descubrí\\
1. Personal & Lo que más cambió en mí fue \linead{6.5cm} & \linead{4.2cm}\\
2. Emocional & La emoción más fuerte fue \linead{2.5cm} y la regulé \linead{3cm} & \linead{4.2cm}\\
3. Ciudadana & En democracia esto importa porque \linead{6cm} & \linead{4.2cm}\\
4. Local (Cartago) & En mi barrio sería útil para \linead{6.5cm} & \linead{4.2cm}\\
5. Intergeneracional & A mi abuela/abuelo le diría \linead{6.5cm} & \linead{4.2cm}\\
\end{tabularx}
\normalsize

\begin{infoband}{milcVino}
\textbf{Para la próxima sustentación.} Vuelve a esta tabla en seis meses. Verás que las respuestas cambian — no porque lo aprendido cambie, sino porque tú cambias. La evaluación liberadora es una conversación contigo a través del tiempo.
\end{infoband}


\puente{Y para terminar, tres voces sobre los datos. Dussel sobre los datos invisibilizados, Marco Aurelio sobre la phronesis del registro, Floridi sobre el dato como acto informacional.}

\begin{titlebox}{milcGradeDark}
Triángulo de pensamiento · cierre filosófico
\end{titlebox}

\begin{softbox}{milcVino}{milcGris}{Dussel · lente del nosotros}
\textit{``Lo que no se cuenta no existe en la decisión pública --- contar es ya un acto político.''}

En Colombia, durante décadas, los pueblos indígenas no fueron ``contados'' en censos oficiales con criterios propios. El conteo del Estado los volvía invisibles. Hoy, en redes sociales, los algoritmos cuentan lo que les conviene contar de tus comportamientos. La pregunta es siempre la misma: \emph{¿quién decide qué se cuenta y para qué?}. Asumir tú la pregunta antes que ellos es ejercicio de soberanía digital.

\textbf{¿Qué de mi vida está siendo contado por apps sin que yo decida? ¿Y qué decido yo contar de mí mismo con criterios propios?}
\end{softbox}
\linead{\linewidth}\\[1mm]
\linead{\linewidth}

\begin{softbox}{milcMostaza}{milcGris}{Estoicismo · Marco Aurelio · cuidado interior}
\textit{``La phronesis empieza por la pregunta antes que por la respuesta.''}

Recolectar datos sin pregunta es la versión moderna del que junta cosas en su casa ``por si acaso''. La sabiduría estoica recomienda lo contrario: antes de recoger, pregúntate qué decisión quieres mejorar. Sin esa disciplina previa, la recolección se vuelve acumulación, y la acumulación se vuelve carga, no virtud.

\textbf{¿Estoy recolectando datos para decidir mejor o para sentirme productivo? ¿Cuál fue la última decisión real que tomé desde una tabla mía?}
\end{softbox}
\linead{\linewidth}\\[1mm]
\linead{\linewidth}

\begin{softbox}{milcTurquesa}{milcGris}{Floridi · lente de la infoesfera}
\textit{``Cada dato registrado es un acto informacional que modifica la infoesfera personal y colectiva.''}

No existe ``recolectar datos neutralmente''. Cada dato que registras hace existir algo que antes no estaba en tu universo de decisión. Cada dato que NO registras deja una zona ciega. Aprender a recolectar con phronesis es aprender a diseñar tu propia infoesfera --- no padecerla.

\textbf{¿Mi recolección de hoy aporta a una infoesfera más útil para mí y mi comunidad, o solo agrega ruido a la mía y a la de las plataformas?}
\end{softbox}
\linead{\linewidth}\\[1mm]
\linead{\linewidth}

\begin{titlebox}{milcVino}
Compromiso final
\end{titlebox}

\small
\begin{tabularx}{\textwidth}{>{\RaggedRight\arraybackslash}p{3.0cm}YYY}
\rowcolor{milcVino}\color{white}\bfseries Criterio & \color{white}\bfseries Lo logré & \color{white}\bfseries En proceso & \color{white}\bfseries Necesito apoyo\\
Comprensión del tema & Explico ideas centrales con mis palabras. & Reconozco ideas pero falta precisión. & Necesito repasar conceptos base.\\
Pensamiento computacional & Apliqué los 4 pilares al diseñar mi producto. & Apliqué algunos pilares. & Necesito releer la fase 2.\\
Aplicación y producto & Mi producto cumple los criterios y se puede revisar. & Producto funcional con ajustes pendientes. & Aún en construcción.\\
Reflexión 5 dimensiones & Respondí cada dimensión con honestidad. & Respondí algunas con detalle. & Mis respuestas son muy generales.\\
Triángulo de pensamiento & Conecté las 3 voces con mi trabajo real. & Conecté al menos una. & No relaciono las voces con mi proceso.\\
\end{tabularx}
\normalsize

\textbf{Hoy entendí que:}\\[1mm]
\linead{\linewidth}\\[1mm]
\linead{\linewidth}

\textbf{Mi compromiso para la próxima sesión es:}\\[1mm]
\linead{\linewidth}\\[1mm]
\linead{\linewidth}

\begin{softbox}{milcMostaza}{milcGris}{Cita de despedida · Marco Aurelio}
\textit{``El obstáculo en el camino se vuelve el camino. Lo que estorba al camino se vuelve camino.''}\\[1mm]
Cada error de hoy, bien observado, se convierte en pieza del camino que pisarás mañana.
\end{softbox}

\small
\begin{tabularx}{\textwidth}{>{\bfseries}p{3.6cm}Y}
\rowcolor{milcGradeDark!12}Nombre completo: & \linead{10cm}\\
Fecha: & \linead{4cm} \quad Curso/grupo: \linead{4cm}\\
Mi firma: & \linead{9cm}\\
\end{tabularx}
\normalsize

\begin{infoband}{milcGradeDark}
\textbf{Para la próxima semana.} Sigue recolectando 5 datos diarios sobre la pregunta que elegiste hoy --- llega a 50 datos reales para la próxima sesión. La tabla bien hecha hoy es la base de todo el periodo. El próximo lunes traemos los datos para empezar a estructurarlos en tablas serias.
\end{infoband}

\end{document}
